\documentclass[10pt,a4paper]{amsart}
\usepackage[margin=19mm]{geometry}
\usepackage{amsmath,amssymb,mathtools}
\usepackage[scr=rsfs,cal=euler]{mathalfa}
\usepackage{xfrac}
\usepackage[unicode,hidelinks,bookmarks=false]{hyperref}
\newcommand{\ie}{i.e.\ }
\newcommand{\cf}{cf.\ }
\newcommand{\resp}{respectively\ }
\newcommand{\Subsection}{\subsection}
\providecommand{\nobreakdash}{\nobreak\hbox{-}}
\setlength{\emergencystretch}{2em}
\multlinegap0pt
\usepackage{bm}
\usepackage{bbm}
\usepackage{dsfont}
\usepackage{xspace}
\usepackage{cancel}
\usepackage{mathtools}
\usepackage{cleveref}
\usepackage{amsthm}
\makeatletter
\@ifundefined{corollary}{\newtheorem{corollary}{Corollary}}{}
\@ifundefined{definition}{\newtheorem{definition}{Definition}}{}
\@ifundefined{lemma}{\newtheorem{lemma}{Lemma}}{}
\@ifundefined{theorem}{\newtheorem{theorem}{Theorem}}{}
\makeatother
\title[A unified framework for establishing the universal approxima]{A unified framework for establishing the universal approximation of transformer-type architectures (Lemma 1)}
\author{Jingpu Cheng, Ting Lin, Zuowei Shen, Qianxiao Li}
\date{Source: arXiv:2506.23551v2 (CC BY 4.0)}
\begin{document}
\maketitle

\noindent\emph{Source and licence.} A unified framework for establishing the universal approximation of transformer-type architectures. Version arXiv:2506.23551v2 (CC BY 4.0), distributed under the Creative Commons Attribution 4.0 International licence, as recorded in the arXiv licence metadata for this exact version and shown on the abstract page https://arxiv.org/abs/2506.23551v2. Full rights notice: licenses/arxiv-2506.23551-CC-BY-4.0.txt.\par\smallskip

\noindent\textit{Editorial scope.} Counted unit: the Lemma printed as Lemma 1 of the article (source0.tex lines 1324--1334, source label \texttt{lem:equal\_atten}), delivered together with the article's own complete original proof of it (source0.tex lines 1336--1402, one \texttt{proof} environment of 67 lines). No mathematics is written, repaired or re-derived: statement and proof are the article's own text line for line. Required context retained uncounted: def:G-UAP (lines 320--326); def:nonlinearity (lines 402--409); def:token\_distinguishability (lines 418--434); thm:main (lines 446--450); thm:kernel (lines 602--624). Every definition, display and label of those lines is retained unchanged. Omitted from this package and not redistributed: 1--319, 327--401, 410--417, 435--445, 451--601, 625--1323, 1335--1335, 1403--1960. Nothing inside a retained excerpt is omitted or altered: every retained body is delivered exactly as the article's lines are written, so no omission receipt of any kind accompanies this package. Citations: the retained excerpts carry no citation command at all, so no bibliography entry of the article is transcribed and no reference is redistributed. Reference adaptation: none was needed and none was made; every reference inside the delivered unit resolves inside it, so no unresolved reference or citation is produced. Printed slips kept literal: no printed word, symbol, hypothesis or number of the counted statement or of its proof was altered. Layout and class: the article's own class is \texttt{article} with packages including neurips\_2025, inputenc, fontenc, hyperref, url, booktabs, amsfonts, nicefrac, microtype, xcolor, graphicx, amsmath, amssymb, amsthm; the wrapper is the lane's pinned \texttt{amsart} editorial wrapper at 10pt on A4, which restates the theorem-like environments the retained text needs and adds the article's own macro definitions that the retained text needs (none), with extra wrapper packages (bm, bbm, dsfont, xspace, cancel, mathtools, cleveref). The delivered numbering is the unit's own. Assets: no retained span contains a figure environment, a graphics-inclusion command, a PDF-page inclusion, a table or a TikZ picture; no image file is copied into this package, the package declares no asset dependency and no image file is redistributed with it. Licence: version arXiv:2506.23551v2 of this article is distributed under the Creative Commons Attribution 4.0 International licence, as recorded in the arXiv licence metadata for this exact version and shown on the abstract page https://arxiv.org/abs/2506.23551v2. A full rights notice is delivered at licenses/arxiv-2506.23551-CC-BY-4.0.txt. Characters: every retained excerpt is pure ASCII, so the delivered engine, XeLaTeX with its default Latin Modern text font, reports no missing character.
\begin{definition}[$G$-UAP]\footnote{Notice that in~\Cref{def:G-UAP}, we do not require $\mathcal T_{\mathcal G,\mathcal H}$ to consist of only $G$-equivariant maps, but only that it can approximate $G$-equivariant functions. }
    \label{def:G-UAP}
    The transformer-type model with hypothesis space $\mathcal{T}_{\mathcal{G},\mathcal{H}}$ is said to have the \emph{$G$-universal approximation property (G-UAP)} in the $L^p$ sense ($1\le p<\infty$) if, for every continuous $G$-equivariant function $F:\mathbb R^{d\times n}\to\mathbb{R}^{d\times n}$, every compact set $K\subset\mathbb R^{d\times n}$, and every $\varepsilon>0$, there exists $\hat{F}\in\mathcal{T}_{\mathcal{G},\mathcal{H}}$ such that
    \begin{equation}
        \|\hat{F}-F\|_{L^p(K)}<\varepsilon.
    \end{equation}
\end{definition}

\begin{definition}[Nonlinearity and affine-invariance for $\mathcal H$]
    \label{def:nonlinearity}
    We say a function family $\mathcal H$ (consisting of functions from $\mathbb{R}^d$ to $\mathbb{R}^d$) is nonlinear and affine-invariant, if 
    \begin{itemize}
        \item For any $h\in \mathcal{H}$ and any $W, A\in \mathbb{R}^{d\times d}, b\in \mathbb{R}^d$, the function $Wh(A\cdot -b)$ also belongs to ${\mathcal{H}}$;
        \item  $\mathcal{H}$ contains at least one non-affine Lipschitz function.
    \end{itemize}
\end{definition}

\begin{definition}[Token distinguishability for $\mathcal G$]
    \label{def:token_distinguishability}
    For a given group $G\le S_n$ and a set of samples $D:=\{X_i\}_{i=1}^N\subset \mathbb{R}^{d\times n}$ that are all in general position, we say a token-mixing family $\mathcal G$ can distinguish tokens in $D$ using $m$ layers under $G$-action, if there exists
    \begin{equation}
        g \in (\operatorname{Id}+\mathcal{G})^m = \{(\operatorname{Id}+g_m)\circ\cdots \circ (\operatorname{Id}+g_1) \mid g_1,\dots, g_m\in \mathcal{G} \}
     \end{equation}
     such that for any distinct $i,j\in[N]$ with $X_i$ and $X_j$ belonging to different orbits under the $G$-action (i.e., $X_i\neq \sigma(X_j)$ for all $\sigma\in G$), the tokens of $g(X_i)$ and $g(X_j)$ are all distinct.

     Moreover, we say $\mathcal G$ satisfies the \textbf{token distinguishability condition} under $G$-action, if for any finite set $D$, there exists $m$ such that $\mathcal G$ can distinguish tokens in $D$ using $m$ layers under $G$-action.

    % We say a token-mixing family $\mathcal G$, consisting of functions from $\mathbb R^{d\times n}$ to $\mathbb R^{d\times n}$, has token distinguishability with $G$-action if for any $N\in\mathbb{N}$ and any set of $N$ samples $\{X_i\}_{i=1}^N\subset \mathbb{R}^{d\times n}$ that are all in general position, there exist $m\in\mathbb{N}$ and
    % \begin{equation}
    %    g \in (\operatorname{Id}+\mathcal{G})^m = \{(\operatorname{Id}+g_m)\circ\cdots \circ (\operatorname{Id}+g_1) \mid g_1,\dots, g_m\in \mathcal{G} \}
    % \end{equation}
    % such that for any distinct $i,j\in[N]$ with $X_i$ and $X_j$ belonging to different orbits under the $G$-action (i.e., $X_i\neq \sigma(X_j)$ for all $\sigma\in G$), the tokens of $g(X_i)$ and $g(X_j)$ are all distinct.
    % That is, the token mixing layers can distinguish tokens across different samples.
\end{definition}

\begin{theorem}
    \label{thm:main}
    Suppose that $\mathcal H$ is nonlinear and affine-invariant~\Cref{def:nonlinearity}, and $\mathcal G$ satisfies the token distinguishability condition~\Cref{def:token_distinguishability}. Then, the family of transformers $\mathcal{T}_{\mathcal{G}, \mathcal{H}}$ satisfies the $G$-UAP~\Cref{def:G-UAP}.
    % \ql{hyperref the definitions for both parts}
\end{theorem}

\begin{corollary}
    \label{thm:kernel}
    % \ql{call this and similar ones corollaries?}
    Suppose the kernel function $k$ satisfies the following conditions:
    \begin{itemize}
        \item $k(\cdot,\cdot):\mathbb R^{d}\times \mathbb R^{d}\to \mathbb R^+$ is an analytic function.

        \item For any $x\in\mathbb{R}^d\setminus\{0\}$ and distinct points $y_1, y_2\in\mathbb{R}^d\setminus\{0\}$, for  almost all $W_K\in\mathbb{R}^{d\times d}$\footnote{means that the condition holds all the whole space except for a measure zero set.}, the following holds:
        \begin{equation}
 \lim_{t\to \infty}\frac{k(x, tW_Ky_1)}{k(x, tW_Ky_2)}= 0  \text{ or } +\infty.
        \end{equation}
        That is, for almost all given $W_K$, the kernel function $k$ can distinguish token representations by scaling the key vectors with a large factor.
        % For any $x\in\mathbb{R}^d\setminus\{0\}$ and distinct points $y_1, y_2\in\mathbb{R}^d\setminus\{0\}$, for almost all(in the Lebesgue measure sense) $W_K\in\mathbb{R}^{d\times d}$:
        % \begin{equation}
        %     \lim_{t\to \infty} \frac{1}{\sum_{j=1}^n k(x, tW_Ky_j)} \left(k(x, tW_Ky_1),\dots, k(x, tW_Ky_n)\right) = \delta_i, \quad \text{for some } i\in[n],
        % \end{equation}
        % where
        % $\delta_i= (0, 0, \dots, 1, \dots, 0, 0)\in\mathbb R^n$ is the $i$-th standard basis vector in $\mathbb{R}^n$. That is, scaling the key vectors $W_Ky_i$ by a large factor $t$ makes the attention scores concentrate on a certain token $i$.

    \end{itemize}
    Then, a transformer with kernel-based attention family $\mathcal{G}$ and feedforward family $\mathcal{H}$ satisfying the conditions in \Cref{thm:main} possesses the ${S}_n$-UAP. Moreover, using only one token-mixing layer and sufficiently many feedforward layers
    can achieve the UAP.
\end{corollary}

\begin{lemma}
    \label{lem:equal_atten}
    Let $r$ be a kernel function satisfying the conditions in \Cref{thm:kernel}.
    Suppose $\{a_1, \dots, a_r\}\subset \mathbb{R}^d\setminus \{0\}$ and $\{b_1,\dots, b_s\}\subset \mathbb{R}^d\setminus \{0\}$ are two sequences of distinct tokens. Suppose that for some indices $r^\prime\in [r]$ and $s^\prime\in [s]$, the following equality holds for all choices of $W_Q, W_K, W_V$:
    \begin{equation}
        \label{eq:lem_equal_atten}
        a_{r^\prime}+\sum_{j=1}^r\left(\frac{ k(W_Q a_{r^\prime}, W_K a_j)}{\sum_{l=1}^r k(W_Q a_{r^\prime}, W_K a_l)} W_V a_j\right) =
        b_{s^\prime}+\sum_{j=1}^s\left(\frac{ k(W_Q b_{s^\prime}, W_K b_j)}{\sum_{l=1}^s k(W_Q b_{s^\prime}, W_K b_l)} W_V b_j\right).
    \end{equation}
    Then, $r=s$, and there exists a permutation $\sigma\in{S}_n$ with $\sigma(r^\prime)=s^\prime$ such that $a_i=b_{\sigma(i)}$ for all $i$.
\end{lemma}

\begin{proof}
    Taking $W_V$ to be zero directly gives $a_{r^\prime}=b_{s^\prime}$. In the following, we set $W_Q=I$, and for notational simplicity define
    \[
    \tilde k(x):=k(a_{r^\prime},x),\quad \forall\,x\in\mathbb{R}^d.
    \]
     The proof utilizes the following lemma:
\begin{lemma}[Auxiliary lemma]
    \label{lem:lemma_aux}
    Let $\{x_1,\cdots, x_p\}:=\{a_1,\cdots, a_r\}\cup \{b_1,\cdots, b_s\}$. Then, there exist $W_K\in\mathbb{R}^{d\times d}$ and a permutation $\sigma\in S_p$ such that
    \begin{equation}
        \lim_{t\to\infty} \frac{\tilde k(tW_K x_{\sigma(j)})}{\tilde k(tW_K x_{\sigma(i)})}=\infty, \text{ for all } i<j.
    \end{equation}
\end{lemma}

Since $\{a_j\}_{j=1}^r$ are $r$ distinct tokens, by the auxiliary lemma, we may choose a $W_K$ (and reindex the sequences accordingly) so that after replacing $W_K$ by $tW_K$ the kernel values satisfy, for large $t>0$,
\[
\tilde k(tW_Ka_1)\ll \tilde k(tW_Ka_2)\ll\cdots\ll \tilde k(tW_Ka_r),
\]
and similarly
\[
\tilde k(tW_Kb_1)\ll \tilde k(tW_Kb_2)\ll\cdots\ll \tilde k(tW_Kb_s).
\]
Moreover, the lemma tells us that for any $a_j$ and $b_l$, either $a_j=b_l$ or one of $\tilde k(tW_ka_j)$ and $\tilde k(tW_kb_l)$ is dominated by the other in the limit $t\to\infty$.


After subtracting $a_{r^\prime}$ from both sides in~\eqref{eq:lem_equal_atten}, we have
\begin{equation}
    \label{eq:mean_eq}
    \frac{\sum_{j=1}^r \tilde k(tW_Ka_j)a_j}{\sum_{j=1}^r \tilde k(tW_Ka_j)}
    = \frac{\sum_{j=1}^s \tilde k(tW_Kb_j)b_j}{\sum_{j=1}^s \tilde k(tW_Kb_j)}.
\end{equation}
By a transformation, it gives
\begin{equation}
    \label{eq:mean_eq2}
\sum_{j=1}^r\sum_{l=1}^s \tilde k(tW_Ka_j)\tilde k(tW_Kb_l)(a_j-b_l)=0.
\end{equation}

Let $t\to \infty$, considering the dominate term $\tilde k(tW_Ka_j)\tilde k(tW_Kb_l)(a_r-b_s)$ gives $a_r=b_s$. Then, for all $q=0, 1,\cdots, \min\{r,s\}$, we then prove by induction that: $a_{r-k}=b_{s-k}$.

Suppose we already have $a_{r-i}=b_{s-i}$, for $i=1,\cdots, q-1$. It then follows that

\begin{equation}
        \label{eq:lem_sum_exp_zero}
        \begin{aligned}
                        &\sum_{j=r-q+1}^r\sum_{l=s-q+1}^s \tilde k(tW_Ka_j)\tilde k(tW_Kb_l)(a_j-b_l)\\
                        &=\sum_{j=r-q+1}^r\sum_{l=r-q+1}^r \tilde k(tW_Ka_j)\tilde k(tW_Ka_l)(a_j-a_l)=0, \quad \text{for all }t\in\mathbb R.
        \end{aligned}
\end{equation}.

Combining with~\eqref{eq:mean_eq2}, we have
\begin{equation}
    (\sum_{j=1}^{r-q}\sum_{l=1}^{r-q}+\sum_{j=1}^{r-q}\sum_{l=s-q+1}^s+\sum_{j=r-q+1}^{r}\sum_{l=1}^{s-q}) \left(\tilde k(tW_Ka_j)\tilde k(tW_Kb_l)(a_j-b_l)\right)=0,
\end{equation}
where the leading term is
\begin{equation}
    \tilde k(tW_Ka_{r-q})\tilde k(tW_Kb_{s})(a_{r-q}-b_{s})+\tilde k(tW_Ka_{r})\tilde k(tW_Kb_{s-q})(a_{r}-b_{s-q}).
\end{equation}
Since $a_r=b_s$, $a_{r-q}\neq a_r$ and $b_{r-q}\neq b_r$, let $t\to\infty$ gives that
$\tilde k(tW_Ka_{r-q})$ and $\tilde k(tW_Kb_{s-q})$ are not dominated by each other. By the auxiliary lemma, this indicates that $a_{r-q}=b_{s-q}$, which completes the induction.

Then, we have shown that $a_{r-k}=b_{s-k}$ for all $k=0, 1,\cdots, \min\{r,s\}$. The only remaining thing is to show that $r=s$. Suppose $r<s$, then we have
\begin{equation}
    \sum_{j=1}^r\sum_{l=1}^{s-r} \tilde k(tW_Ka_j)\tilde k(tW_Kb_l)(a_j-b_l)\equiv 0,
\end{equation}
where the unique leading term is $\tilde k(tW_Ka_r)\tilde k(tW_Kb_{s-r})(a_r-b_{s-r})$. Since $a_r=b_r\neq b_{s-r}$, that gives a contradiction.
This completes the proof.
\end{proof}

\end{document}
